% =====================================================================
% paper_content.tex -- body of the paper.
%
% The file is split into PART blocks. Each block is owned by one part
% branch, which keeps parallel writing conflict-free:
%
%   PART front : background, restatement, assumptions, notation
%                -> branch part/hNN/front
%   PART p1    : model, solution and results of Problem 1
%                -> branch part/hNN/p1
%   PART p2    : Problem 2                                    -> part/hNN/p2
%   PART p3    : Problem 3                                    -> part/hNN/p3
%   PART sens  : sensitivity analysis                         -> part/hNN/sens
%   PART eval  : model evaluation, strengths/weaknesses, conclusion
%                -> branch part/hNN/eval (optional)
%
% Merge order into the hypothesis branch (keeps conflicts low):
%   front -> p1 -> p2 -> p3 -> sens -> eval -> misc (abstract, title, ...)
%
% Rules:
%   1. Only edit the block you own; never reformat or reorder other blocks.
%   2. Keep every result traceable: a number here must come from code
%      under code/ and a figure under figures/ (see docs/latex.md).
%   3. If a merge conflicts inside your block, keep your own version and
%      take the other blocks unchanged.
% =====================================================================

% ===== PART front =====
\section{Introduction and Problem Restatement}

Describe the background and restate each sub-problem in your own words.
Cite related work as needed, for example \citep{hyndman2018forecasting}
and \citep{lecun2015deep}. Replace this placeholder text.

\section{Model Assumptions}

List the assumptions of the model, one item per assumption, each with a
one-sentence justification.

\section{Notation}

Define the symbols that are reused in the following sections.

% ===== PART p1 =====
\section{Problem 1: Model and Solution}

\subsection{Model Formulation}

\subsection{Solution and Results}

Present the model, the algorithm and the results of Problem 1 here.

% ===== PART p2 =====
\section{Problem 2: Model and Solution}

\subsection{Model Formulation}

\subsection{Solution and Results}

Present the model, the algorithm and the results of Problem 2 here.

% ===== PART p3 =====
\section{Problem 3: Model and Solution}

\subsection{Model Formulation}

Present the model and the algorithm of Problem 3 here. If the model is an
optimisation problem, state the objective and constraints explicitly and
cite the solution method, for example \citep{bertsekas1999nonlinear}.

\subsection{Solution and Results}

Present the model, the algorithm and the results of Problem 3 here.

% ===== PART sens =====
\section{Sensitivity Analysis}

Analyse how the results change with the key parameters, data and
assumptions; report the ranges in which the conclusions stay valid.

% ===== PART eval =====
\section{Model Evaluation and Conclusion}

Summarise strengths, weaknesses and possible extensions of the model, and
give the final conclusion of the paper.
